problem:  
Let \(S^{n+1}\) be the unit round sphere with its standard metric \(g_{\mathrm{round}}\). For each smooth metric \(g\) sufficiently close to \(g_{\mathrm{round}}\) in the \(C^{3}\) topology, denote by \(\mathcal{M}(g)\) the collection of integral \(n\)-currents in \(S^{n+1}\) that are *area-minimizing boundaries* (that is, currents of the form \(\partial E\) minimizing mass among all boundaries homologous to \(\partial E\)).  

It is known that for the round metric, certain singular minimizing cones (for example, the Simons cone in suitable local coordinate regions) can occur as tangent cones to minimizers. One also knows that in higher dimensions (\(n \ge 7\)), singularities of area-minimizing hypersurfaces cannot, in general, be avoided.  

The question is:  

> **Does there exist an open dense subset \(\mathcal{G}\) of smooth metrics on \(S^{n+1}\)** (in the \(C^{3}\) topology) such that for every \(g \in \mathcal{G}\), *no area-minimizing boundary in \((S^{n+1}, g)\) admits a singular tangent cone that is linearly stable in the Euclidean sense*?  
>   
> Equivalently, do generic metrics on the sphere preclude the appearance of locally stable singular minimizing cones as blow-up limits of area-minimizing hypersurfaces?

In other words, rather than asking for full smoothness generically, this problem seeks to determine whether **stable singular cones**, which are known to persist in flat space, can actually occur as local models of area-minimizers for a generic Riemannian metric on the sphere.  

Why is it a "good" problem:  
This problem refines the previously overly strong generic smoothness question into one that is both plausible and deep within current geometric measure theory.

1. **Conceptual depth and precision:**  
   By focusing on the *stability* of singular tangent cones rather than on complete regularity, it asks whether the singular phenomena of minimizing hypersurfaces are “structurally stable” or an artifact of the metric’s symmetry. This directly tests the robustness of the Almgren–Simons regularity theory under perturbations.

2. **Bridge between fields:**  
   It unites geometric measure theory (stability of minimizing cones and tangent cone analysis) with global differential geometry (geometry of metrics on manifolds and generic perturbations). Tools from elliptic PDEs, metric perturbation theory, and possibly infinite-dimensional Sard–Smale transversality would be required to make progress.

3. **Extreme simplicity and clarity:**  
   The question can be stated in plain geometric terms: *Are singular minimizing cones stable under generic perturbations of the ambient metric on the sphere?* Despite its simplicity, it pertains to one of the deepest mechanisms behind geometric singularities.

4. **Future potential:**  
   A positive (or negative) resolution would profoundly influence the study of regularity and generic behavior in geometric variational problems. It could reveal a sharp frontier between “inevitable” and “structurable” singularities, influencing not only minimal surface theory but also mean curvature flow and calibrated geometry.

In summary, this refined problem is mathematically sound, aligns with current frontiers in minimal surface theory, and calls for innovative cross-disciplinary techniques, fulfilling the essential criteria for a “good” research-level problem.